\chapter{Modelo de datos}
\label{cap:datos}

Una base PostgreSQL 18, dos esquemas, dos dueños. \pkcode{auth} tiene cinco
tablas (\pkcode{user}, \pkcode{session}, \pkcode{account}, \pkcode{verification},
\pkcode{jwks}) y las escribe Better Auth. \pkcode{pokedex} tiene dieciocho tablas
de dominio más \pkcode{alembic\_version}, y las escribe Alembic. Los modelos
viven en \pkpath{apps/api/src/pokedex/db/models/}.

\section{Dos esquemas, dos herramientas de migración}
\label{sec:migraciones}

La separación no es estética: cada esquema lo define el sistema que sabe cómo
tiene que ser. Better Auth genera sus tablas desde su propia configuración —si
mañana se añade un plugin, la tabla que hace falta la crea él— y Alembic genera
las del dominio desde \pkcode{Base.metadata}. El precio de mezclarlos sería que
\pkcode{alembic revision --autogenerate} viera tablas ausentes de su metadata y
propusiera borrarlas.

Ese riesgo se cierra con dos filtros en \pkpath{apps/api/alembic/env.py}, que
cubren direcciones distintas:

\begin{lstlisting}[language=Python]
DOMAIN_SCHEMA = "pokedex"

def include_name(name, type_, parent_names):
    # Without this filter autogenerate reflects Better Auth's `auth` tables and
    # proposes dropping them, since they are absent from Base.metadata.
    if type_ == "schema":
        return name == DOMAIN_SCHEMA
    return True

def include_object(obj, name, type_, reflected, compare_to):
    # `auth.user` is registered in Base.metadata only so cross-schema foreign
    # keys resolve; alembic must never try to create or drop it.
    if type_ == "table":
        return bool(getattr(obj, "schema", None) == DOMAIN_SCHEMA)
    return True
\end{lstlisting}

\pkcode{include\_name} bloquea lo que se \emph{refleja} desde la base;
\pkcode{include\_object} bloquea lo que viene de la metadata. La tabla
\pkcode{auth.user} está declarada en \pkpath{db/models/external.py} como un
\pkcode{Table} suelto con dos columnas, únicamente para que resuelvan las claves
foráneas entre esquemas.

Dos cosas no pueden crearse dentro de la cadena de migraciones, porque la
primera revisión ya las necesita, y por eso se ejecutan antes:

\begin{lstlisting}[language=Python]
await connection.execute(text(f'CREATE SCHEMA IF NOT EXISTS "{DOMAIN_SCHEMA}"'))
await connection.execute(text("CREATE EXTENSION IF NOT EXISTS pg_trgm"))
\end{lstlisting}

\begin{pkwarn}[El \pkcode{search\_path} está fijado a \pkcode{public}]
El motor se crea con \pkcode{server\_settings=\{"search\_path": "public"\}}. Con
el valor por defecto (\pkcode{"\$user", public}) el esquema \pkcode{pokedex}
pasaría a ser implícito y los filtros de Alembic dejarían de discriminar. Por
la misma razón \pkcode{pg\_trgm} se instala en \pkcode{public}: un
\pkcode{gin\_trgm\_ops} instalado en otro esquema sería invisible al crear el
índice.
\end{pkwarn}

\section{Catálogo y colección}

\begin{figure}[p]
  \centering
  \includegraphics[width=\linewidth,height=0.80\textheight,keepaspectratio]{er-catalogo.pdf}
  \caption{Capa de catálogo: una carta es una impresión, una especie es lo
    retratado, y \pkcode{card\_price} es la serie detrás del último precio.
    Generado desde \pkpath{diagramas/src/er-catalogo.dot}.}
  \label{fig:er-catalogo}
\end{figure}

\begin{figure}[p]
  \centering
  \includegraphics[width=\linewidth,height=0.80\textheight,keepaspectratio]{er-coleccion.pdf}
  \caption{Capa de colección: todo cuelga de \pkcode{auth."user"} y del catálogo,
    y nada de aquí es visible para otro usuario. Generado desde
    \pkpath{diagramas/src/er-coleccion.dot}.}
  \label{fig:er-coleccion}
\end{figure}

\subsection{Por qué \pkcode{card\_price} es una fila por carta y por día}

\pkcode{card.price\_usd} guarda \emph{siempre} la última lectura, y eso hace el
movimiento inmedible: un portafolio sin serie es un número suelto. La tabla
\pkcode{card\_price} es la historia detrás de ese número, con clave primaria
compuesta \pkcode{(card\_id, recorded\_on)}. La escritura es un
\pkcode{ON CONFLICT DO UPDATE} sobre esa misma clave, así que volver a
sincronizar el mismo día \emph{sobrescribe} el día en vez de añadir un punto:
la serie se mantiene en un punto por día y el trabajo es idempotente.

\subsection{Por qué un sobrante es \pkcode{SUM(quantity) > 1}}

No existe una columna «para cambiar». Un sobrante se calcula:

\begin{lstlisting}[language=SQL]
SELECT user_id, card_id, SUM(quantity) - 1 AS spare
FROM pokedex.collection_item
GROUP BY user_id, card_id
HAVING SUM(quantity) > 1
\end{lstlisting}

La primera copia es la colección; sólo lo que pasa de ahí es inventario. La
consecuencia práctica es que nadie tiene que mantener un estado
«disponible/reservado» coherente con las cantidades reales: al publicar, al
ofertar y al aceptar se recalcula, así que el tablón nunca puede prometer una
carta que ya no existe. El agrupamiento es por \pkcode{(user\_id, card\_id)} y no
por condición: dos copias en estados distintos siguen siendo dos copias.

\subsection{Por qué el único de \pkcode{collection\_item} usa \pkcode{NULLS NOT DISTINCT}}

Una fila no es una carta física: es un grupo homogéneo de copias. El único
cubre \pkcode{(user\_id, card\_id, condition, language, is\_graded, grade)}, y
\pkcode{grade} es \pkcode{NULL} para todo lo que no está gradeado. Postgres
trata los \pkcode{NULL} como distintos por omisión, lo que permitiría insertar
el mismo grupo una y otra vez e impediría que \pkcode{ON CONFLICT} disparase
nunca. \pkcode{postgresql\_nulls\_not\_distinct=True} es lo que convierte el alta
en una fusión real.

El coste se guarda por grupo, no por copia (\pkcode{unit\_cost\_usd}), de modo
que copias compradas a precios distintos se promedian —\pkcode{\_blended\_cost}
pondera por cantidad— en lugar de quedar como lotes separados. Si un lado no
tiene coste conocido, el coste conocido carga el grupo entero.

\section{Intercambio, conversación y mensajería}

\begin{figure}[p]
  \centering
  \includegraphics[width=\linewidth,height=0.80\textheight,keepaspectratio]{er-intercambio.pdf}
  \caption{Intercambio. Una oferta tiene destinatario y una publicación no;
    tomarla escribe la oferta ya aceptada que las une. Generado desde
    \pkpath{diagramas/src/er-intercambio.dot}.}
  \label{fig:er-intercambio}
\end{figure}

\begin{figure}[p]
  \centering
  \includegraphics[width=\linewidth,height=0.80\textheight,keepaspectratio]{er-mensajeria.pdf}
  \caption{Conversación con el asistente y mensajería entre coleccionistas.
    Generado desde \pkpath{diagramas/src/er-mensajeria.dot}.}
  \label{fig:er-mensajeria}
\end{figure}

\subsection{Una oferta tiene destinatario; una publicación, no}

Esa es toda la diferencia entre \pkcode{trade\_offer} y \pkcode{trade\_listing}.
Tomar una publicación escribe una \pkcode{trade\_offer} \emph{ya aceptada}, con
el publicador como emisor y quien la tomó como destinatario, y enlaza ambas por
\pkcode{trade\_listing.offer\_id}. Desde ese momento, un trato que empezó en el
tablón y uno que empezó como propuesta son lo mismo, y quien la tomó vive en un
solo sitio, así que los dos registros no pueden contradecirse.

\subsection{Por qué la condición se copia y la cantidad no se guarda}

\pkcode{trade\_offer\_card} lleva \pkcode{condition} como columna propia en vez de
apuntar a la fila de colección de la que salió: una oferta es una promesa sobre
una carta \emph{en un estado}, y tiene que sobrevivir a que su dueño edite o
borre esa fila. Y la condición no es un detalle —es la mayor parte del precio:
una carta \emph{near mint} y una \emph{damaged} son la misma fila sin ella, y
valen varias veces la una a la otra.

La cantidad, en cambio, está deliberadamente ausente: una oferta nombra cartas,
y cuántas copias cambian de manos se acuerda cuando las dos personas quedan.
Anotar un número aquí parecería un compromiso que la aplicación no puede exigir
a nadie.

En \pkcode{trade\_listing\_card} sólo el lado \emph{ofrecido} lleva condición: en
qué estado llegará una carta pedida lo responde quien tome la publicación, y
todavía no hay a quién preguntar.

\subsection{Los multiplicadores de condición}

Viven en \pkpath{services/trade.py} y su orden de declaración es además el orden
de mejor a peor que usa el resto del módulo:

\par\addvspace{8pt}\noindent\begin{pkbox}{colspec={X[1.6,l]X[1,r]X[3.4,l]}}
  Condición & Factor & Uso \\
  \pkcode{mint} & 1.05 & \SetCell[r=6]{l} \pkcode{adjusted(price, condition)} multiplica y cuantiza a dos decimales. Se aplica al valorar una oferta, una publicación y el simulador. La lista también define el orden «peor copia», que es el que se toma cuando una oferta no declara condición. \\
  \pkcode{near\_mint} & 1.00 & \\
  \pkcode{lightly\_played} & 0.85 & \\
  \pkcode{moderately\_played} & 0.70 & \\
  \pkcode{heavily\_played} & 0.50 & \\
  \pkcode{damaged} & 0.35 & \\
\end{pkbox}\par\addvspace{8pt}

\begin{pknote}[Son ratios del gremio, no una medición]
El comentario que acompaña a la tabla lo dice sin rodeos: «estas son las
proporciones con las que negocia la afición, no algo medido aquí, así que toda
cifra derivada de ellas se rotula como estimación allí donde se muestra». La
interfaz respeta ese contrato.
\end{pknote}

\subsection{Un hilo es un par ordenado}

\pkcode{direct\_thread} guarda \pkcode{first\_user\_id} y \pkcode{second\_user\_id}
ordenados, con \pkcode{CHECK first\_user\_id < second\_user\_id} y único sobre el
par. Con un par emisor/receptor las mismas dos personas podrían abrir dos hilos
escribiéndose a la vez, y cada lectura tendría que unir las dos mitades. No hay
tabla de participantes porque no hay un tercer participante.

\pkcode{direct\_message.read\_at} no necesita dueño: en un hilo de dos, quien no
escribió el mensaje es la única persona que puede leerlo. Y su
\pkcode{created\_at} usa \pkcode{clock\_timestamp()} y no \pkcode{now()}, porque
\pkcode{now()} es el instante en que empezó la transacción: dos mensajes
escritos en la misma transacción compartirían marca y el hilo se quedaría sin
orden en el que leerse.

\subsection{El historial del asistente se guarda íntegro}

\pkcode{message.content} es \pkcode{jsonb} con el \emph{payload} del framework tal
cual, para que una conversación reanudada se reproduzca con sus llamadas a
herramientas intactas en lugar de como una transcripción con pérdida.
\pkcode{message.text} está desnormalizado para poder listar una conversación sin
deserializar cada \emph{payload}.

\pkcode{user\_preference} es la memoria del agente: una fila por clave, única por
usuario. Las claves con prefijo \pkcode{sys.} son internas y no se exponen por
la API; \pkcode{sys.notifications\_seen\_at} es la que sostiene el estado
«visto» del feed de novedades.

\section{Índices y restricciones que importan}

\par\addvspace{8pt}\noindent\begin{pkbox}{colspec={X[2.2,l]X[1.1,l]X[2.7,l]}}
  Objeto & Tipo & Por qué existe \\
  \pkcode{ix\_card\_name\_normalized\_trgm} & GIN \pkcode{gin\_trgm\_ops} & La resolución de escaneos compara nombres con el operador \pkcode{\%}; sin este índice sería un recorrido completo del catálogo. \\
  \pkcode{ix\_species\_types} & GIN sobre \pkcode{text[]} & Filtrar la colección por tipo. \\
  \pkcode{ix\_message\_conversation\_created} & B-tree compuesto & Una conversación se lee siempre en orden y nunca como feed global. \\
  \pkcode{ix\_direct\_message\_thread\_created} & B-tree compuesto & Lo mismo para un hilo directo. \\
  \pkcode{uq\_collection\_item\_group} & UNIQUE \emph{nulls not distinct} & Convierte el alta en fusión. \\
  \pkcode{ck\_collection\_item\_quantity\_positive} & CHECK & Un grupo con cero copias no es un grupo. \\
  \pkcode{ck\_collection\_item\_grade\_only\_when\_graded} & CHECK & \pkcode{grade} y \pkcode{is\_graded} no pueden discrepar. \\
  \pkcode{ck\_direct\_thread\_ordered} & CHECK & Hace verificable el par ordenado. \\
  \pkcode{uq\_trade\_offer\_card} & UNIQUE & La misma carta no puede repetirse en el mismo lado de una oferta. \\
\end{pkbox}\par\addvspace{8pt}

\section{Migraciones}

Catorce revisiones en \pkpath{apps/api/alembic/versions/}, encadenadas desde
\pkcode{bb4e34c86509\_init}. La tabla de versión vive en el propio esquema
(\pkcode{version\_table\_schema="pokedex"}) y \pkcode{compare\_type} está activo,
así que un cambio de tipo se detecta en el autogenerado.

\begin{lstlisting}[language=bash]
uv run alembic upgrade head           # aplicar
uv run alembic revision --autogenerate -m "mensaje en ingles"
uv run alembic downgrade -1           # revertir una
\end{lstlisting}

\begin{pkwarn}[Orden de arranque]
Las tablas de dominio tienen claves foráneas hacia \pkcode{auth."user"}, que crea
Better Auth. En un despliegue nuevo los dos servicios arrancan a la vez, así que
la imagen de la API reintenta \pkcode{alembic upgrade head} hasta treinta veces
con cinco segundos de espera en lugar de entrar en bucle de caída.
\end{pkwarn}
